%! AP Statistics — Unit 2, Lesson 2.6 — Guided Notes
\documentclass[10pt]{article}
\usepackage{apstats-article}
\usepackage{apstats-boxes}

\begin{document}

\pageheader{Unit 2, Lesson 2.6}{Guided Notes}
\namedateperiod

\vspace{0.12in}

% ── OBJECTIVE ────────────────────────────────────────────────────────────────
\begin{objectivebox}
\small By the end of this lesson, I will be able to\dots\\[4pt]
\writeline\\[1pt]
\writeline\\[1pt]
\writeline
\end{objectivebox}

\vspace{0.1in}

% ── VOCABULARY ───────────────────────────────────────────────────────────────
\begin{vocabbox}
\termblanklong{Regression line (LSRL)}
\termblanklong{Slope ($b$)}
\termblanklong{$y$-intercept ($a$)}
\termblanklong{$\hat{y}$ (``$y$-hat'' / predicted value)}
\termblanklong{Interpolation}
\termblanklong{Extrapolation}
\end{vocabbox}

\vspace{0.1in}

% ── HOOK ─────────────────────────────────────────────────────────────────────
\begin{hookbox}
\small\textbf{Opening Question:} Your teacher shows a scatterplot of height (inches) vs.\
shoe size for 30 students and asks: ``If a student is 68 inches tall, what shoe size would
you predict?''

\vspace{6pt}
Write your estimate: Predicted shoe size $\approx$ \blank{2.5cm} \quad How did you decide?\\[8pt]
\writeline\\[1pt]\writeline

\vspace{4pt}
\textbf{Key idea:} Today we formalize this process with an equation --- the \textbf{regression line}.
\end{hookbox}

\vspace{0.1in}

% ── STEP 1 ───────────────────────────────────────────────────────────────────
\begin{notesbox}{Step 1 --- The Regression Equation (DAT-1.E)}
\small

\noindent The equation of the \textbf{least-squares regression line (LSRL)} is:
\[
  \hat{y} = a + bx
\]
where \quad $a = $ \blank{5.5cm} \quad and \quad $b = $ \blank{5.5cm}.

\vspace{6pt}
\textbf{AP rule:} Always write the equation using the \blank{5cm} of the variables,
not generic $x$ and $y$.

\vspace{8pt}
\textbf{Class example} --- age of a used car ($x$, years) vs.\ resale price ($\hat{y}$, \$1000s):

\vspace{4pt}
Technology output gives $a = 28.4$ and $b = -1.85$. Write the equation using variable names:

\vspace{6pt}
\hspace{1.5cm} \blank{9cm}

\end{notesbox}

\vspace{0.1in}

% ── STEP 2 ───────────────────────────────────────────────────────────────────
\begin{notesbox}{Step 2 --- Interpreting the Slope (DAT-1.E)}
\small

\noindent The \textbf{slope} $b$ is the predicted \blank{4cm} in $\hat{y}$ for each
\blank{3cm} increase in $x$.

\vspace{6pt}
\textbf{AP template:} ``For each additional \blank{5cm}, the predicted
\blank{4.5cm} [increases / decreases] by approximately \blank{3cm}.''

\vspace{8pt}
\textbf{Interpret the slope} ($b = -1.85$) for the car example:
\vspace{4pt}\\
\writeline\\[2pt]\writeline

\vspace{6pt}
\textbf{Every AP slope interpretation needs all three:}
\begin{enumerate}[leftmargin=1.8em, itemsep=2pt]
  \item \blank{9cm} \hfill (one-unit change in $x$, in context)
  \item \blank{9cm} \hfill (direction: increases or decreases)
  \item \blank{9cm} \hfill (magnitude with units)
\end{enumerate}

\end{notesbox}

\vspace{0.1in}

% ── STEP 3 ───────────────────────────────────────────────────────────────────
\begin{notesbox}{Step 3 --- Interpreting the $y$-Intercept (DAT-1.E)}
\small

\noindent The \textbf{$y$-intercept} $a$ is the predicted value of $\hat{y}$ when $x = $ \blank{1.5cm}.

\vspace{6pt}
\textbf{AP template:} ``When [x-variable] $= 0$, the predicted [y-variable] is approximately
\blank{3cm}.''

\vspace{8pt}
\textbf{Interpret the $y$-intercept} ($a = 28.4$) for the car example:
\vspace{4pt}\\
\writeline\\[2pt]\writeline

\vspace{6pt}
\textbf{Is this meaningful in context?} \blank{3cm} \quad \textbf{Why or why not?}\\[6pt]
\writeline\\[2pt]\writeline

\vspace{4pt}
\begin{tcolorbox}[colback=redbg, colframe=redacc, boxrule=0.8pt, arc=2mm,
    left=3mm, right=3mm, top=2mm, bottom=2mm]
\small\textbf{Watch out:} If $x = 0$ is \blank{4cm} the data range or physically
\blank{4cm}, state that the $y$-intercept is \textbf{not meaningful in context}.
\end{tcolorbox}

\end{notesbox}

\vspace{0.1in}

% ── STEP 4 ───────────────────────────────────────────────────────────────────
\begin{notesbox}{Step 4 --- Predictions: Interpolation vs.\ Extrapolation (DAT-1.F)}
\small

\noindent To \textbf{predict} $\hat{y}$, substitute an $x$ value into the regression equation.

\vspace{6pt}
\textbf{Predict the price of a 5-year-old car:}

\vspace{4pt}
\hspace{1.5cm} $\widehat{\text{price}} = 28.4 - 1.85(\underline{\hspace{1.2cm}})
  = \blank{2.5cm}$ thousand dollars

\vspace{8pt}
\renewcommand{\arraystretch}{1.7}
\begin{tabularx}{\linewidth}{@{} >{\bfseries}p{3.2cm} X @{}}
\rowcolor{sky} \textbf{Term} & \textbf{Definition} \\
Interpolation &
  $x$ is \blank{4cm} the data range; prediction is generally \blank{2.5cm}. \\
Extrapolation &
  $x$ is \blank{4cm} the data range; prediction is \blank{4cm} and should be flagged. \\
\end{tabularx}

\vspace{8pt}
\textbf{The car data ranged from age 1 to 10 years.}

\vspace{4pt}
Predict the price of a \textbf{20-year-old} car:
\hspace{0.5cm} $\widehat{\text{price}} = 28.4 - 1.85(20) = \blank{3cm}$ thousand dollars

\vspace{4pt}
Is this interpolation or extrapolation? \blank{3.5cm} \quad Why is this prediction unreliable?\\[6pt]
\writeline\\[2pt]\writeline

\end{notesbox}

\vspace{0.1in}

% ── STEP 5 ───────────────────────────────────────────────────────────────────
\begin{notesbox}{Step 5 --- The LSRL Passes Through $(\bar{x},\, \bar{y})$}
\small

\noindent The regression line always passes through the point \blank{4cm} --- the mean of $x$
and the mean of $y$.

\vspace{6pt}
\textbf{Verify for the car data:} $\bar{x} = 5.5$ years, \; $\bar{y} = \blank{2.5cm}$ thousand dollars.

\vspace{4pt}
\hspace{1.5cm} $\widehat{\text{price}} = 28.4 - 1.85(5.5) = \blank{3cm}$ \quad Matches $\bar{y}$? \blank{2cm}

\end{notesbox}

\vspace{0.1in}

% ── CONNECTIONS ───────────────────────────────────────────────────────────────
\begin{spiralbox}
\small
\begin{multicols}{2}

\textbf{How this connects:}
\begin{itemize}[itemsep=3pt]
  \item Lesson 2.5: $r$ describes direction and strength; slope $b$ gives the rate of change.
  \item Lesson 2.7: each point's \textit{residual} $= y - \hat{y}$ measures prediction error.
  \item Lesson 2.8: the LSRL minimizes the sum of squared residuals.
\end{itemize}

\columnbreak

\textbf{Reflection:}\\
Why is extrapolation unreliable? Give your own example of a prediction that would involve
extrapolation.\\[2pt]
\writeline\\[1pt]\writeline\\[1pt]\writeline

\end{multicols}
\end{spiralbox}

\vspace{0.1in}

\begin{notesbox}{Additional Notes}
\vspace{0pt}
\writeline\\\writeline\\\writeline\\\writeline\\\writeline\\\writeline
\end{notesbox}

\end{document}
